% !PRINTED 67
\weberpage{67}
Bestimmt man nach 5. eine Funktion $T(u)$ mit den Nullpunkten $\alpha_1, \alpha_2, \ldots, \alpha_m$, so ist:
\[
T(u)\varphi(u) = T_1(u)
\]
eine $T$-Funktion $m$ter Ordnung, und
\[
\varphi(u) = \frac{T_1(u)}{T(u)} \qquad \text{w.\,z.\,b.\,w.}
\]
Die Funktionen $T(u)$, $T_1(u)$ sind \gesperrt{verwandt}, da sie dieselben Periodizitätsfaktoren und also auch dieselbe Charakteristik haben.

Die Nullpunkte von $T_1(u)$ sind zugleich die Nullpunkte von $\varphi(u)$, und die Summe dieser \gesperrt{Nullwerte} ist also kongruent mit \gesperrt{der Summe der Unstetigkeitswerte}.

\begin{center}
\textbf{\S\,19. $T$-Funktionen erster Ordnung.}
\end{center}

Wir gehen nun dazu über, die $T$-Funktionen erster Ordnung, die wir mit $t(u)$ bezeichnen, aus denen sich die übrigen $T$-Funktionen, wie wir gesehen haben, ableiten lassen, näher zu bestimmen.

Aus 3. des vorigen Paragraphen folgt, daß zwei $t$-Funktionen derselben Charakteristik sich nur durch einen Faktor von der Form
\[
Ce^{-\pi i(\lambda u^2+\mu u)}
\]
voneinander unterscheiden, und wir wollen nun durch eine Erweiterung der Definition diesen Faktor noch näher bestimmen.

Dies soll, was nach § 17, (9) ohne Beschränkung der Allgemeinheit möglich ist, so geschehen, daß in den Periodizitätsfaktoren
\[
a_1 = 0, \quad b_1 = g_1, \quad b_2 = g_2
\]
wird, wenn $(g_1, g_2)$ die Charakteristik ist; wegen der Relation
\[
a_2\omega_1 - a_1\omega_2 = 1
\]
ist dann
\[
a_2 = \frac{1}{\omega_1},
\]
und die Bedingungen für diese $t$-Funktion lauten alsdann (§ 17, 1.):
\weq{1}{\begin{aligned}
t(u+\omega_1) &= e^{-\pi i g_1}t(u),\\
t(u+\omega_2) &= e^{-\pi i \frac{2u+\omega_2}{\omega_1}}e^{-\pi i g_2}t(u),
\end{aligned}}
und hierdurch ist die Funktion $t(u)$ bis auf einen von $u$ unabhängigen Faktor bestimmt. Um diesen Faktor noch näher zu
%CONT>>
